% page 579 du fac-simile (image p-578 du scan)
% bloc 152.4 x 263.6 mm ; marge gauche 49.6 mm
\pgc{27.940mm}{0.000mm}{
\l{}
\l{}
\l{}
\l{\cel{55}571}
\l{}
\l{\sou{tarso}. (\sou{anat.}) Parto dopa di la ped-osto, qua konsistas ye}
\l{\cel{3}sep osti. - EFIS.}
\l{}
\l{\sou{tarto}. Kuko, ronda (maxim-multa-kaze), ek fundo plata de}
\l{\cel{3}pasto, kun rebordo, e kovrita ye frukti, konfitaji,}
\l{\cel{3}kremo, e c. - DEFIRS.}
\l{}
\l{\sou{tartano-navo}. (\sou{navig.})\cel{2}Mikra navo (sur Mediteraneo) pont-}
\l{\cel{3}oza, un-masta, e kun seglo triangula (kun un ortangulo).}
\l{}
\l{\sou{.tartano-stofo}. Stofo ek lano, kun figuri ek quadrato-frami,}
\l{\cel{3}di qui la kolori alternas. -}
\l{}
\l{\sou{tarta}ro. (\sou{kemio}). I. Sedimento sala qua diferenciesas po-}
\l{\cel{3}kope ek la liquori vinatra,e adheras a la parieti di la}
\l{\cel{3}bareli, forme di krusto. - II.Sedimento blankatra o}
\l{\cel{3}flavatra qua amaseskas an la infra parto di la denti e}
\l{\cel{3}karias pokope lia emalio. - EFIS.}
\l{}
\l{\sou{taso}. Mikra vazo kun anso, ek porcelano, metalo, vitro,}
\l{\cel{3}quan onu uzas por drinkar (kafeo, teo, chokolado, e c.)}
\l{\cel{3}- DFIS.}
\l{}
\l{\sou{tasko}. I. To quon onu obligesas laborar, facar, agar.- II.}
\l{\cel{3}Quanto de laboro quan onu su obligis exekutor en fristo}
\l{\cel{3}e po preco konvencionita. - III. Laboro skribala quan}
\l{\cel{3}docisto donas kom exekutenda da lua docarii exter la doc-}
\l{\cel{3}eyo. - EF.}
\l{}
\l{\sou{tastar}. (\sou{netrans}.) I. Palpar per la manui, en tenebro, por}
\l{\cel{3}avancar, por trovar to quon onu serchas. - II. Probar}
\l{\cel{3}interdiferante, pro ne savar quale agar apte por atingar}
\l{\cel{3}la skopo quan onu determinis. - DFI.}
\l{}
\l{\sou{tatuar}. (\sou{trans.}) Trasar (sur la pelo) desegnuri, per farbi}
\l{\cel{3}quin onu koaktas - per la pinto di instrumento - penetrar}
\l{\cel{3}sub la epidermo. - DEFIR.}
\l{}
\l{\sou{taumaturgo}. Facisto di mirakli. (\sou{vorto pejorativa}).- DEFIS.}
\l{}
\l{\sou{taumaturgio.}\cel{2}La arto di la taumaturgo. - DEFIS.}
\l{}
\l{\sou{taverno}.I.(\sou{olim)}. Butiko en qua onu iris konsumar drinkaji}
\l{\cel{3}plu o min alkoholoza, manjar, fumar (+smokar).-\cel{2}II.}
\l{\cel{3}(\sou{nun}) Kafe-drinkerio, restorerio plu o min luxoza, ed}
\l{\cel{3}(en ula kazi) ornita da objekti artala (freski, pikturi,}
\l{\cel{3}statui). - EFIS.}
\l{}
\l{\sou{taxar}. (\sou{trans., ye)}\cel{3}Enskribar (ulu) sur la registro de}
\l{\cel{3}la impostebli kom debanta pagor pekunio-quanto (de...)}
\l{\cel{3}- DEFR.}
\l{}
\l{\sou{taximetro}. (\sou{automobilo}). Kontilo specala, uzata sur la fia-}
\l{\cel{3}kri, e qua indikas la pekunio-quanto pagenda, kande la}
\l{\cel{3}transporto jus finabos, proporcione ye la quanto de kolo-}
\l{\cel{3}metri parkurita. - DEFIS.}
}
